\chapter[Compiler-Scheduled ILP]{Compiler-Scheduled ILP: VLIW and Software Pipelining}\label{ch:compiler-scheduled-ilp}

\section{VLIW as a Compiler-Scheduled Alternative}\label{sec:ilp-vliw}
One response to dynamic-scheduling complexity is to move more of the scheduling work into the compiler.
As superscalar windows grow, the machine would like a larger scheduling queue so it can search farther ahead for independent work, but larger queues also make wakeup, select, and operand-readiness checks harder to close at a short cycle time.
This leads to the \textit{very long instruction word (VLIW)} approach, historically associated with Josh Fisher.
A VLIW encoding packs several operations intended to execute in parallel into one instruction word, so the compiler decides the schedule at compile time rather than letting large runtime scheduling hardware discover it dynamically.

In this context, an individual encoded action is an \textit{op}.
A set of ops intended to execute in the same cycle is a \textit{multi-op}.
VLIW reduces runtime hardware complexity because the machine no longer needs a large associative scheduling queue, aggressive wakeup logic, or much of the dynamic issue machinery used in FOCO-style processors.
The tradeoff is that the compiler must do much more work up front, and the resulting code is less adaptive to runtime variation such as cache misses or changing branch behavior.
This tradeoff is attractive in domains such as \textit{digital signal processors (DSPs)}, where workloads are regular and low-power execution is important.

Compiler design must balance three competing goals: code quality, compilation time, and maintainability.
In practice, that usually rules out algorithms much worse than quadratic for large compiler passes, and even cubic methods are used cautiously.
In principle, finding the optimal schedule for a dependence-constrained block is exponential in block size, so practical VLIW compilers rely on heuristics rather than on exact search.
List scheduling is one of the standard heuristics for this job.
The broader performance point is that widening fetch and scheduling machinery yields diminishing returns: useful IPC usually grows sublinearly with front-end width, while the cost of dynamic scheduling hardware grows quickly.


\section{Dependence Graphs and List Scheduling}\label{sec:ilp-vliw-list-scheduling}
The compiler begins from a \textit{data dependence graph}.
True flow edges capture read-after-write dependencies, while anti-dependence and output-dependence edges capture false constraints caused by register reuse.
Pure flow dependence always takes priority over false dependence because it represents real data movement.
If the code contains loops, this graph can contain cycles, so it is not necessarily a directed acyclic graph.
Constructing the graph by hand is mechanical.
First draw one node per op.
Then add all true flow edges from each definition to the later uses that require that value.
Next add anti-dependence edges from a read to a later write of the same architectural name, and output-dependence edges between two writes of the same name.
If two operations are already ordered by a true flow dependence, that true edge is the one that matters for scheduling because it already enforces the required order.

Table~\ref{tab:ilp-vliw-block} gives a compact block that will be used for the scheduling example.
Assume one pipelined load/store unit, one pipelined arithmetic logic unit (ALU), and one pipelined multiplier.
Each unit can accept a new op every cycle; the listed latencies describe when results become available, not how long the unit remains occupied.
Also assume the compiler has already assigned registers so that the remaining constraints shown in the table are dominated by true dependences rather than by avoidable name conflicts.
\begin{table}[t]
    \centering
    \renewcommand{\arraystretch}{1.15}
    \begin{tabular}{|c|l|l|c|}
        \hline
        Instr.
        & Operation & True dependences & Latency \\
        \hline
        A & LW $R2, X$ & None & 2 \\
        \hline
        B & LW $R3, Y$ & None & 2 \\
        \hline
        C & LW $R4, Z$ & None & 2 \\
        \hline
        D & ADD $R5, R1, R2$ & A & 1 \\
        \hline
        E & MUL $R6, R2, R3$ & A, B & 3 \\
        \hline
        F & MUL $R7, R1, R3$ & B & 3 \\
        \hline
        G & ADD $R8, R7, \#1$ & F & 1 \\
        \hline
        H & SW $R8, X$ & G & 1 \\
        \hline
        I & ADD $R9, R2, R4$ & A, C & 1 \\
        \hline
    \end{tabular}
    \caption{Illustrative basic block for VLIW list scheduling.}
    \label{tab:ilp-vliw-block}
\end{table}
Figure~\ref{fig:ilp-vliw-ddg} presents the same block in a more visual form: the original operations, the resulting data dependence graph, and one priority list ordered by dependence height.
\begin{figure}[t]
    \centering
    \small
    \begin{minipage}[t]{0.27\textwidth}
        \raggedright
        \textbf{Original operations}\par
        \vspace{0.5ex}
        \begin{tabular}[t]{@{}ll@{}}
            A: & LW $R2, X$ \\
            B: & LW $R3, Y$ \\
            C: & LW $R4, Z$ \\
            D: & ADD $R5, R1, R2$ \\
            E: & MUL $R6, R2, R3$ \\
            F: & MUL $R7, R1, R3$ \\
            G: & ADD $R8, R7, \#1$ \\
            H: & SW $R8, X$ \\
            I: & ADD $R9, R2, R4$
        \end{tabular}
    \end{minipage}\hfill
    \begin{minipage}[t]{0.45\textwidth}
        \centering
        \textbf{Data Dependence Graph}\par
        \vspace{1ex}
        \begin{tikzpicture}[
            >=Latex,
            depnode/.style={draw, circle, inner sep=1pt, minimum size=15pt, font=\small},
            hlabel/.style={font=\scriptsize, text=blue!65!black}
        ]
            \node[depnode] (A) at (0,0) {A};
            \node[depnode] (B) at (1.7,0) {B};
            \node[depnode] (C) at (3.4,0) {C};
            \node[depnode] (D) at (0,-1.1) {D};
            \node[depnode] (E) at (1.7,-1.1) {E};
            \node[depnode] (F) at (0,-2.2) {F};
            \node[depnode] (I) at (2.7,-2.2) {I};
            \node[depnode] (G) at (0,-3.3) {G};
            \node[depnode] (H) at (0,-4.4) {H};

            \node[hlabel, left=1.5mm of A] {$5$};
            \node[hlabel, above=1.5mm of B] {$7$};
            \node[hlabel, above=1.5mm of C] {$3$};
            \node[hlabel, left=1.5mm of D] {$1$};
            \node[hlabel, right=1.5mm of E] {$3$};
            \node[hlabel, left=1.5mm of F] {$5$};
            \node[hlabel, right=1.5mm of I] {$1$};
            \node[hlabel, left=1.5mm of G] {$2$};
            \node[hlabel, left=1.5mm of H] {$1$};

            \draw[->] (A) -- (D);
            \draw[->] (A) -- (E);
            \draw[->] (A) -- (I);
            \draw[->] (B) -- (E);
            \draw[->] (B) -- (F);
            \draw[->] (C) -- (I);
            \draw[->] (F) -- (G);
            \draw[->] (G) -- (H);
        \end{tikzpicture}
    \end{minipage}\hfill
    \begin{minipage}[t]{0.20\textwidth}
        \raggedright
        \textbf{List ordered by dependence height}\par
        \vspace{0.5ex}
        \begin{tabular}[t]{@{}l@{}}
            $B$ $(7)$ \\
            $A$ $(5)$ \\
            $F$ $(5)$ \\
            $C$ $(3)$ \\
            $E$ $(3)$ \\
            $G$ $(2)$ \\
            $D$ $(1)$ \\
            $I$ $(1)$ \\
            $H$ $(1)$
        \end{tabular}
    \end{minipage}
    \caption{Original operations, dependence graph, and one priority ordering for the VLIW list-scheduling example. The small numeric annotations next to nodes indicate dependence height. Ties may be broken arbitrarily.}
    \label{fig:ilp-vliw-ddg}
\end{figure}

One common priority heuristic is \textit{dependence height}, which estimates how much latency remains below a node on the longest dependent path.
If the heuristic includes the latency of the current instruction itself, one convenient definition is Eq.~\eqref{eq:ilp-dependence-height}.
\begin{equation}\label{eq:ilp-dependence-height}
    height(i)=
    \begin{cases}
        lat(i), & succ(i)=\emptyset\\
        lat(i)+\max\limits_{j\in succ(i)} height(j), & \text{otherwise}
    \end{cases}
\end{equation}
Higher-height instructions receive higher priority because delaying them is more likely to lengthen the critical path.
To annotate a graph by hand, start from the sinks and assign each sink its own latency.
Then sweep backward, adding the current node latency to the maximum height among its successors.
The ready list is obtained by sorting nodes from highest height to lowest, breaking ties either arbitrarily or by an explicit secondary rule such as program order.
For the block in Table~\ref{tab:ilp-vliw-block}, the longest true-dependence chain is $B\rightarrow F\rightarrow G\rightarrow H$, so the heuristic tends to prioritize $F$ ahead of $E$ once both become candidates for the single multiplier.

List scheduling then tries to place the highest-priority ready instructions into each cycle subject to both dependence timing and resource availability.
A simple legality condition is shown in Eq.~\eqref{eq:ilp-list-schedule-legal}.
\begin{equation}\label{eq:ilp-list-schedule-legal}
    cycle(i)=\min\!\left\{c\ \middle|\ \forall j\in pred(i),\ c\ge cycle(j)+lat(j),\ \text{and a compatible FU is free in }c \right\}.
\end{equation}
This is not an optimal algorithm, but it is practical and captures the central VLIW idea: schedule as much parallel work as possible before runtime while still respecting dependencies and resource limits.

Table~\ref{tab:ilp-vliw-schedule} shows one possible compiler-produced schedule for the block running on a machine with one arithmetic logic unit (ALU), one multiplier, and one load/store unit.
The point is not that this schedule is unique, but that the compiler can decide the cycle packing in advance and encode it directly.
For illustration, this schedule keeps the three initial loads in source order even though a greedier dependence-height policy could start $B$ before $A$ and shorten the block by one cycle.
That choice leaves one explicit empty cycle later, which makes the pause-field encoding example concrete.
\begin{table}[t]
    \centering
    \renewcommand{\arraystretch}{1.15}
    \begin{tabular}{|c|c|c|c|}
        \hline
        Cycle & Load/Store Unit & ALU & Multiplier \\
        \hline
        1 & A &  &  \\
        \hline
        2 & B &  &  \\
        \hline
        3 & C & D &  \\
        \hline
        4 &  &  & F \\
        \hline
        5 &  & I & E \\
        \hline
        6 &  &  &  \\
        \hline
        7 &  & G &  \\
        \hline
        8 & H &  &  \\
        \hline
    \end{tabular}
    \caption{One legal VLIW-style schedule for the block under fixed functional-unit limits.}
    \label{tab:ilp-vliw-schedule}
\end{table}
The schedule shows two important points.
First, the compiler only packs operations together when both the dependence constraints and the functional-unit constraints allow it.
Second, the absence of an operation in a slot does not require the hardware to discover what to do; the compiler can encode the schedule directly.
In particular, cycle 6 is intentionally empty because $G$ cannot start until the three-cycle multiply $F$ has finished producing $R7$.
Table~\ref{tab:ilp-vliw-packed-schedule} shows the same schedule in a field-level packing format closer to an actual multi-op layout.
This table defines which operands occupy each unit's fields in each cycle; it does not by itself define the serialized op order in memory.
\begin{table}[t]
    \centering
    \small
    \setlength{\tabcolsep}{4pt}
    \renewcommand{\arraystretch}{1.12}
    \begin{tabular}{|c|c c c|c c c|c c|c c|}
        \hline
        Cycle & \multicolumn{3}{c|}{IALU (1 cycle)} & \multicolumn{3}{c|}{MUL (3 cycles)} & \multicolumn{2}{c|}{Load (2 cycles)} & \multicolumn{2}{c|}{Store (1 cycle)} \\
        \hline
         & $d$ & $s_1$ & $s_2$ & $d$ & $s_1$ & $s_2$ & $d$ & addr & addr & $s$ \\
        \hline
        1 &  &  &  &  &  &  & $R2$ & $X$ &  &  \\
        \hline
        2 &  &  &  &  &  &  & $R3$ & $Y$ &  &  \\
        \hline
        3 & $R5$ & $R1$ & $R2$ &  &  &  & $R4$ & $Z$ &  &  \\
        \hline
        4 &  &  &  & $R7$ & $R1$ & $R3$ &  &  &  &  \\
        \hline
        5 & $R9$ & $R2$ & $R4$ & $R6$ & $R2$ & $R3$ &  &  &  &  \\
        \hline
        6 &  &  &  &  &  &  &  &  &  &  \\
        \hline
        7 & $R8$ & $R7$ & $\#1$ &  &  &  &  &  &  &  \\
        \hline
        8 &  &  &  &  &  &  &  &  & $X$ & $R8$ \\
        \hline
    \end{tabular}
    \caption{Field-level packing of the schedule from Table~\ref{tab:ilp-vliw-schedule}. Rows 3 and 5 are multi-ops because they contain more than one op in the same cycle.}
    \label{tab:ilp-vliw-packed-schedule}
\end{table}


\section{Multi-Op Encoding, Tail Bits, and Pause Fields}\label{sec:ilp-vliw-encoding}
VLIW raises an encoding question: how should a machine represent the boundaries of a multi-op if not every cycle uses every functional unit?
One simple approach is to fill every unused slot with an explicit no-op, but that wastes code space.
An alternative is to encode only real ops and add a small amount of grouping information.
Each encoded op still carries its ordinary fields such as opcode, destination register, and source registers; the additional VLIW-specific fields only describe grouping and delay.

One useful scheme adds a \textit{tail bit} to each encoded op.
If the tail bit is 0, the current op is not the end of the current multi-op.
If the tail bit is 1, the current op ends the current multi-op.
The next op then begins a new multi-op.
This lets the instruction stream represent a variable number of ops per cycle without requiring every cycle to contain the same number of explicit no-op fillers.

A related field is a \textit{pause field}, which tells the front-end control logic to wait a specified number of cycles before beginning the next multi-op.
In the simple scheme used here, the pause value is interpreted only on the op whose tail bit ends the current multi-op; pause values on earlier ops in the same multi-op are ignored.
This is a compact way to represent silence between multi-ops while still keeping each op's encoding regular.
Table~\ref{tab:ilp-vliw-encoding} illustrates the idea using the schedule from Table~\ref{tab:ilp-vliw-schedule}.
It shows one legal serialized op order for the same multi-ops that appear in Table~\ref{tab:ilp-vliw-packed-schedule}.
\begin{table}[t]
    \centering
    \renewcommand{\arraystretch}{1.15}
    \begin{tabular}{|c|c|c|p{2.9in}|}
        \hline
        Encoded op & Tail & Pause & Meaning \\
        \hline
        A & 1 & 0 & Single-op multi-op for cycle 1. \\
        \hline
        B & 1 & 0 & Single-op multi-op for cycle 2. \\
        \hline
        C & 0 & 0 & First op of the cycle-3 multi-op. \\
        \hline
        D & 1 & 0 & Ends the cycle-3 multi-op. \\
        \hline
        F & 1 & 0 & Single-op multi-op for cycle 4. \\
        \hline
        I & 0 & 0 & First op of the cycle-5 multi-op. \\
        \hline
        E & 1 & 1 & Ends the cycle-5 multi-op and inserts one pause cycle before the next multi-op. \\
        \hline
        G & 1 & 0 & Single-op multi-op for cycle 7. \\
        \hline
        H & 1 & 0 & Single-op multi-op for cycle 8. \\
        \hline
    \end{tabular}
    \caption{Illustrative VLIW encoding using tail bits and a pause field rather than explicit no-op padding.}
    \label{tab:ilp-vliw-encoding}
\end{table}
Some implementations instead use simple interlocking at the multi-op level, which can remove the need for an explicit pause field.
That saves encoding space, but it also gives up some of the compiler's direct control over exactly when the next multi-op may begin.


\section{Software Pipelining and Cyclic Scheduling}\label{sec:ilp-software-pipelining}
Loops create another opportunity for compiler-scheduled parallelism.
Instead of scheduling only one copy of a loop body, the compiler can overlap operations from different iterations so that a later iteration begins before an earlier one has fully finished.
The older name for this idea is \textit{polycyclic scheduling}; the more common modern name is \textit{software pipelining}.

Conceptually, the compiler behaves as if it had unrolled the loop many times, compacted the overlapping schedule, and then kept only the repeating pattern.
That repeating pattern is the \textit{kernel}.
The code that fills the pipeline before the kernel reaches steady state is the \textit{prologue}.
The code that drains the remaining work after loop exit is the \textit{epilogue}.
The distance in cycles between the start of successive iterations is the \textit{initiation interval (II)}.
Once the schedule reaches steady state, the loop begins a new iteration every $II$ cycles, so the throughput is $1/II$ iterations per cycle.


\section{Iterative Modulo Scheduling}\label{sec:ilp-iterative-modulo-scheduling}
A practical way to construct such schedules is \textit{Iterative Modulo Scheduling (IMS)}.
The compiler first estimates a lower bound on $II$, then attempts to build a modulo schedule at that value.
If the attempt fails, it increments $II$ and tries again.

The first bound comes from resource pressure.
If a loop contains $N_t$ operations of type $t$ and the machine provides $U_t$ compatible units of that type, then
\begin{equation}\label{eq:ilp-res-ii}
    ResII=\max_t \left\lceil \frac{N_t}{U_t} \right\rceil
\end{equation}
This is the resource-constrained initiation interval.

The second bound comes from loop-carried recurrences.
If a dependence cycle $c$ has total latency $Lat(c)$ and spans $Dist(c)$ loop iterations, then
\begin{equation}\label{eq:ilp-rec-ii}
    RecII=\max_c \left\lceil \frac{Lat(c)}{Dist(c)} \right\rceil
\end{equation}
This is the recurrence-constrained initiation interval.
If the dependence graph has no loop-carried cycles, then the practical lower bound is $RecII=1$.

The schedule must satisfy both limits, so the compiler begins from
\begin{equation}\label{eq:ilp-min-ii}
    MII=\max(ResII,RecII),\qquad II\ge MII.
\end{equation}
Many texts call this lower bound the \textit{minimum initiation interval (MII)}.
IMS therefore follows a simple pattern:
\begin{enumerate}
    \item Start with $II=MII$.
    \item Attempt a legal modulo schedule at that value.
    \item If the attempt fails, increment $II$ and retry.
\end{enumerate}
A modulo reservation table for a chosen $II$ has exactly $II$ rows.
An operation that would occupy cycle $x$ in a longer unrolled schedule is placed into row $x \bmod II$, provided the dependence and resource constraints remain legal.

In most textbook treatments, the underlying machine is modeled with a point-latency abstraction: an op occupies one issue slot, then its result becomes ready after a fixed latency.
This keeps the scheduling problem manageable.
More detailed models that explicitly schedule register-file ports, result buses, and intermediate pipeline stages are possible, but they are much more expensive to optimize and are usually avoided in introductory treatments.


\section{Worked Software-Pipelining Example}\label{sec:ilp-software-pipelining-example}
Table~\ref{tab:ilp-swpipe-loop} shows a small loop that illustrates the idea.
Assume two arithmetic logic units, two memory units, one-cycle add/branch latency, and two-cycle load latency.
The branch still logically closes the loop body even though later iterations will overlap older ones.
\begin{table}[t]
    \centering
    \renewcommand{\arraystretch}{1.15}
    \begin{tabular}{|c|l|l|c|}
        \hline
        Instr.
        & Operation & Key dependences & Latency \\
        \hline
        A & ADD $R2, R2, \#1$ & loop-carried on $R2$ & 1 \\
        \hline
        B & LW $R3, 0(R2)$ & A & 2 \\
        \hline
        C & SW $R0, 0(R3)$ & B & 1 \\
        \hline
        D & ADD $R10, R10, \#1$ & loop-carried on $R10$ & 1 \\
        \hline
        E & BRGTZ $R10$, LOOP & D; branch must close the loop body & 1 \\
        \hline
    \end{tabular}
    \caption{Loop used to illustrate software pipelining and modulo scheduling.}
    \label{tab:ilp-swpipe-loop}
\end{table}

The cyclic dependence view of this loop contains the within-iteration edges $A_i\rightarrow B_i$, $B_i\rightarrow C_i$, and $D_i\rightarrow E_i$, plus the loop-carried edges $A_i\rightarrow A_{i+1}$ and $D_i\rightarrow D_{i+1}$.
Those loop-carried edges are exactly the recurrences used when computing $RecII$.

For this loop, the resource bound is
\[
ResII=\max\!\left(\left\lceil \frac{3\ \text{ALU/branch ops}}{2\ \text{ALUs}} \right\rceil,
\left\lceil \frac{2\ \text{memory ops}}{2\ \text{memory units}} \right\rceil \right)=\max(2,1)=2.
\]
The loop-carried recurrences on $R2$ and $R10$ each have latency 1 and distance 1, so $RecII=1$.
Therefore the first useful attempt is $II=2$.

One legal modulo reservation table is shown in Table~\ref{tab:ilp-swpipe-kernel}.
The first row begins iteration $i$.
The second row finishes the current iteration's control work while also performing the store from iteration $i-1$.
That overlap is the essence of software pipelining.
\begin{table}[t]
    \centering
    \renewcommand{\arraystretch}{1.15}
    \begin{tabular}{|c|c|c|c|c|}
        \hline
        Modulo slot & ALU 0 & ALU 1 & MEM 0 & MEM 1 \\
        \hline
        0 & $A_i$ & $D_i$ &  &  \\
        \hline
        1 & $E_i$ &  & $B_i$ & $C_{i-1}$ \\
        \hline
    \end{tabular}
    \caption{One modulo reservation table for the loop in Table~\ref{tab:ilp-swpipe-loop} with $II=2$.}
    \label{tab:ilp-swpipe-kernel}
\end{table}

The emitted code must still handle loop entry and loop exit correctly.
Table~\ref{tab:ilp-swpipe-phases} separates the schedule into prologue, kernel, and epilogue form.
The early-exit branch $E^\ast$ is the inverted form of the loop branch; it lets the machine leave the prologue cleanly if the loop should terminate before the steady-state kernel is entered.
\begin{table}[t]
    \centering
    \renewcommand{\arraystretch}{1.15}
    \begin{tabular}{|c|l|}
        \hline
        Phase & Operations \\
        \hline
        Prologue & $A_1, D_1$ \quad then \quad $B_1, E^\ast_1$ \\
        \hline
        Kernel & repeat every $II=2$ cycles: \quad $A_i, D_i$ \quad then \quad $B_i, C_{i-1}, E_i$ \\
        \hline
        Epilogue & $C_{\text{last}}$ \\
        \hline
    \end{tabular}
    \caption{Prologue, kernel, and epilogue for the software-pipelined loop.}
    \label{tab:ilp-swpipe-phases}
\end{table}

This example also shows why software pipelining is attractive for VLIW machines.
The compiler exposes the repeating parallel structure explicitly, so the hardware can execute the steady-state kernel with very little dynamic scheduling logic.
At the same time, general-purpose superscalar machines still benefit from dynamic scheduling because real loads do not always return in a fixed number of cycles and binaries may need to run across several microarchitectural variants.


\section{Memory Disambiguation}\label{sec:ilp-vliw-memory-disambiguation}
Loads and stores create a special scheduling problem.
If a compiler moves a load ahead of an older store to the same address, the load may read the wrong value.
The compiler therefore needs \textit{memory disambiguation}: it must prove that two memory references do not overlap, or else it must preserve the original order or insert explicit synchronization.

This is one of the limits of VLIW compared with dynamic scheduling hardware.
An out-of-order processor may speculate and recover if it guesses wrong.
A VLIW compiler must usually be more conservative unless it has strong alias information.
That is one reason VLIW works best when memory behavior is regular enough for compile-time reasoning to be reliable.

\section{Chapter 7 Summary}\label{sec:chapter7-summary-vliw}
Compiler-scheduled ILP takes a different approach from dynamic superscalar scheduling.
Instead of building increasingly complex runtime issue logic, a VLIW machine relies on the compiler to expose parallel ops ahead of time and encode them explicitly.
Data dependence graphs, dependence-height ordering, and list scheduling provide the basic machinery for building legal multi-ops.
Tail-bit and pause-field encoding then translate that schedule into a compact instruction stream without filling unused slots with explicit no-ops.
Software pipelining extends the same idea across loop iterations by constructing a prologue, steady-state kernel, and epilogue around a chosen initiation interval.
Memory disambiguation remains the key correctness limit: compile-time scheduling can move loads and stores only when aliasing is understood well enough to preserve program behavior.
Together these techniques show how compiler scheduling can approach the performance of aggressive ILP hardware while trading hardware complexity for compile-time analysis.
